\documentclass[11pt,letterpaper]{article}
\usepackage[utf8]{inputenc}
\usepackage[T1]{fontenc}
\usepackage[spanish]{babel}
\usepackage{amsmath,amssymb}
\usepackage[margin=1in]{geometry}
\usepackage{enumitem}
\usepackage{array}
\usepackage{fancyhdr}
\pagestyle{plain}
\setlength{\parindent}{0pt}

\begin{document}

\begin{center}
  {\Large\bfseries Comentes en frío -- PEP 1}\\[2pt]
  {\large\bfseries Macroeconomía II}\\[6pt]
  \rule{\textwidth}{0.4pt}\\[2pt]
  Parte I de la Guía de Ejercicios (Ayudante: Bastian Puntarelli)
\end{center}

\begin{itemize}
  \item Indique si la afirmación es \textbf{Verdadera, Falsa o Incierta}, y justifique con ecuaciones,
  gráficos y conceptos (costo marginal, productividad marginal, efecto nivel versus efecto crecimiento,
  etc.). Una respuesta sin justificación no es una respuesta.
  \item \textbf{Tiempo: 20 minutos} ($\le 60$ s por comente): veredicto + ecuación clave + una frase.
  Si no sale, escriba ``?'' y siga. Luego 10 minutos de corrección.
\end{itemize}

\section*{Parte I: Comentes}
\textit{Verdadero, Falso o Incierto. Justifique su respuesta.}

\subsection*{A. Modelo de Solow básico}
\begin{enumerate}[series=comentes]
  \item Una economía que ahorra más siempre tendrá un mayor consumo por trabajador en el estado estacionario.
  \item Dos economías son idénticas en $s$, $n$, $\delta$ y $\alpha$, pero una de ellas parte con la mitad del capital por
  trabajador de la otra. Ambas tendrán el mismo producto por trabajador en el largo plazo, pero la
  más pobre crecerá más rápido mientras converge.
  \item Si una economía ahorra más que lo que indica la regla de oro, puede aumentar su consumo por
  trabajador tanto hoy como en el largo plazo reduciendo su tasa de ahorro.
  \item En el modelo de Solow sin progreso tecnológico, el producto por trabajador no crece en el largo
  plazo. Por lo tanto, el modelo no puede explicar el crecimiento sostenido del ingreso per cápita.
  \item El modelo de Solow predice que la tasa de crecimiento del capital por trabajador es mayor mientras
  más lejos esté la economía de su estado estacionario.
\end{enumerate}

\subsection*{B. Modelo de Solow con progreso tecnológico}
\begin{enumerate}[resume=comentes]
  \item En el sendero de crecimiento balanceado, el capital y el producto por trabajador crecen a la tasa $g$,
  mientras que el capital y el producto agregados crecen a la tasa $n+g$.
  \item Un aumento permanente de la tasa de ahorro aumenta permanentemente la tasa de crecimiento del
  producto por trabajador.
  \item Un aumento de la tasa de crecimiento tecnológico $g$ reduce el valor de estado estacionario de $k/A$.
  Sin embargo, en el largo plazo el producto por trabajador termina siendo mayor.
  \item Dos países tienen la misma función de producción, $\delta$ y $g$, y sólo difieren en su tasa de crecimiento
  poblacional $n$. El país con mayor $n$ tendrá menor PIB per cápita y también menor tasa de crecimiento
  del PIB per cápita en el largo plazo.
\end{enumerate}

\subsection*{C. Capital humano, residuo de Solow y convergencia}
\begin{enumerate}[resume=comentes]
  \item Un aumento permanente de los años de escolaridad $u$ aumenta permanentemente la tasa de crecimiento del producto por trabajador.
  \item Con $\psi = 0{,}10$, dos años adicionales de escolaridad aumentan el producto por trabajador de estado
  estacionario exactamente en 20\%.
  \item El residuo de Solow mide la parte de las diferencias de producto por trabajador entre países que no
  se explica por diferencias de capital físico ni de capital humano, y se interpreta como diferencias de
  productividad.
  \item El modelo de Solow predice que los países pobres siempre crecen más rápido que los países ricos.
  \item Si $k/(Ah)$ de un país está por debajo de su valor de estado estacionario, entonces su capital por
  trabajador crece a una tasa mayor que $g$.
\end{enumerate}

\subsection*{D. Modelo de Romer}
\begin{enumerate}[resume=comentes]
  \item Que las ideas sean no rivales implica que es imposible excluir a otros de su uso.
  \item En el modelo de Romer, la competencia imperfecta en el sector de bienes intermedios es necesaria
  para que haya crecimiento: si el precio fuera igual al costo marginal, no habría incentivo a innovar.
  \item Si $\phi < 1$ y el número de investigadores $L_A$ fuera constante, la tasa de crecimiento de la tecnología,
  $\dot{A}/A$, tendería a cero.
  \item En el modelo de Romer con $\phi < 1$, un aumento permanente de la fracción de trabajadores dedicados
  a investigar, $s_R$, aumenta permanentemente la tasa de crecimiento de largo plazo de la economía.
  \item Mientras mayor sea la fracción de trabajadores dedicados a investigar, $s_R$, mayor será siempre el
  producto por trabajador.
  \item En el modelo de Romer, el producto agregado se escribe $Y = K^{\alpha}(AL_Y)^{1-\alpha}$, igual que en el modelo
  de Solow con progreso tecnológico. Por lo tanto, el modelo de Romer no agrega nada nuevo.
\end{enumerate}

\newpage
\section*{Hoja de respuestas}
\renewcommand{\arraystretch}{1.9}
\begin{tabular}{|c|c|p{10.6cm}|c|}
  \hline
  \textbf{\#} & \textbf{V/F/I} & \textbf{Ecuación clave + concepto} & \textbf{$\checkmark$/$\times$} \\ \hline
  1 & & & \\ \hline
  2 & & & \\ \hline
  3 & & & \\ \hline
  4 & & & \\ \hline
  5 & & & \\ \hline
  6 & & & \\ \hline
  7 & & & \\ \hline
  8 & & & \\ \hline
  9 & & & \\ \hline
  10 & & & \\ \hline
  11 & & & \\ \hline
  12 & & & \\ \hline
  13 & & & \\ \hline
  14 & & & \\ \hline
  15 & & & \\ \hline
  16 & & & \\ \hline
  17 & & & \\ \hline
  18 & & & \\ \hline
  19 & & & \\ \hline
  20 & & & \\ \hline
\end{tabular}

\bigskip
\textbf{Puntaje: \underline{\hspace{1.5cm}} / 20} \quad (meta $\ge 17$). Un veredicto correcto sin ecuación o concepto cuenta como $\times$.
\end{document}
